\section{Conclusiones y Trabajos Futuros}

Aunque los satélites han demostrado ser herramientas poderosas para la conectividad, su vulnerabilidad sigue recordándonos la importancia de diseñar con resiliencia desde el primer día. Resulta revelador observar cómo una sola falla puede revertir años de esfuerzos, tal como ocurrió con el VENESAT-1.

\subsection{Conclusiones}

Uno de los hallazgos más contundentes de este estudio es que una arquitectura redundante bien diseñada para el VENESAT-2 resulta técnicamente factible y operativamente viable bajo las condiciones tropicales venezolanas. Las simulaciones exhaustivas demostraron que la combinación de redundancia 1+1, diversidad geográfica de estaciones terrenas y protocolos de conmutación rápidos permite alcanzar disponibilidades superiores al 99.95\%, superando significativamente los niveles observados en sistemas sin respaldo.

Los balances de enlace confirman márgenes adecuados incluso en banda Ka, siempre que se implementen técnicas adaptativas y diversidad efectiva. El análisis FMEA reveló reducciones sustanciales en los riesgos críticos, aunque también identificó áreas donde persisten vulnerabilidades correlacionadas. Estos resultados contrastan fuertemente con la situación vivida en 2020, cuando la pérdida de control del satélite Simón Bolívar dejó sin servicio a gran parte del país \cite{VenesatStatus2020}.

Más allá de los números, el estudio subraya que la redundancia no es un lujo técnico sino una necesidad estratégica para garantizar soberanía en comunicaciones. La propuesta logra equilibrar desempeño, complejidad y costos relativos, ofreciendo una hoja de ruta realista. No obstante, cabe matizar que el éxito final dependerá de la ejecución precisa, la capacitación local y el compromiso sostenido con el mantenimiento.

En síntesis, los datos recopilados y analizados apoyan firmemente la implementación del sistema de enlace redundante. Venezuela puede avanzar hacia una infraestructura satelital más robusta, siempre que se respeten las particularidades climáticas, operativas y presupuestarias identificadas.

\subsection{Trabajos futuros}

Lejos de ser un asunto cerrado, este estudio abre múltiples líneas de investigación y desarrollo que merecen atención prioritaria. Particularmente prometedor resulta explorar la integración progresiva de capacidades multi-satélite y elementos de constelaciones LEO como respaldo complementario al VENESAT-2.

Entre las recomendaciones concretas destacan: la realización de campañas de medición in-situ de propagación durante al menos dos años completos para refinar los modelos locales de atenuación; el desarrollo de algoritmos de inteligencia artificial específicos para predicción de conmutación y mitigación de desvanecimientos; y el análisis económico detallado de diferentes escenarios de implementación por fases.

Resulta estratégico también investigar modelos de gobernanza y operación compartida que permitan reducir costos mediante alianzas regionales o público-privadas. La evolución hacia sistemas 6G, con su énfasis en convergencia y multi-conectividad, ofrece un horizonte técnico estimulante que debería guiar las actualizaciones del VENESAT-2 en los próximos ciclos de vida \cite{HTSSurvey2025}.

Adicionalmente, se sugiere profundizar en aspectos de ciberseguridad satelital, sostenibilidad energética de las estaciones terrenas y desarrollo de capacidades locales de fabricación e integración de subsistemas. Estos trabajos no solo fortalecerían el proyecto actual, sino que sentarían bases para una industria satelital venezolana más madura y autónoma.

En última instancia, el VENESAT-2 redundante no debe verse como un punto final, sino como el inicio de una evolución continua. Los desafíos identificados aquí, lejos de desanimar, invitan a mantener un esfuerzo sostenido de investigación aplicada que transforme la factibilidad demostrada en una realidad operativa robusta y perdurable.